\documentclass[11pt]{article}
\usepackage[T1]{fontenc}
\usepackage[utf8]{inputenc}
\usepackage{lmodern}
\usepackage{amsmath,amssymb}
\usepackage{longtable,booktabs,array}
% Compatibility with Pandoc's captionless tables on older longtable versions.
\ifcsname c@none\endcsname\else\newcounter{none}\fi
\usepackage{graphicx}
\usepackage[margin=25mm]{geometry}
\usepackage{hyperref}
\hypersetup{colorlinks=true,linkcolor=blue,urlcolor=blue}
\providecommand{\tightlist}{\setlength{\itemsep}{0pt}\setlength{\parskip}{0pt}}
\providecommand{\pandocbounded}[1]{#1}
\setlength{\emergencystretch}{3em}
\setcounter{secnumdepth}{-1}
\begin{document}
\section{Cohomology of sheaves on ringed
spaces}\label{cohomology-of-sheaves-on-ringed-spaces}

\emph{Written by GPT-6.1 Sol (OpenAI), in Codex, at Ultra effort,
October 2026. Self-checked by the writing AI, GPT-6.1 Sol, at Ultra
effort. Public domain (CC0).}

A sheaf lets us recognize a section by looking locally. It does not
promise that a section of a quotient sheaf lifts globally. Sheaf
cohomology measures this failure and its higher analogues. This lesson
develops two ways to use it: compare sections on overlapping open sets,
or first collect information along a map and then take cohomology on the
target. Both methods work on arbitrary ringed spaces.

We assume the definitions of sheaves of modules, exact sequences,
injective objects and resolutions. The lesson
\href{https://kokunoyumeto.github.io/open-math-courses/courses/derived-categories-and-sheaf-operations/injective-modules-and-bounded-below-derived-functors.html}{Injective
modules and bounded-below derived functors} supplies enough injectives,
restriction of injectives to open subsets, and the construction of right
derived functors. The exact-couple construction and convergence needed
for Leray are proved in
\href{https://kokunoyumeto.github.io/open-math-courses/courses/derived-categories-and-sheaf-operations/derived-pullback-and-pushforward.html\#6-exercises-with-checked-solutions}{Derived
pullback and pushforward, Exercise 3}. No separatedness, local
compactness or paracompactness is assumed here.

Basic references are {[}Stacks{]} and {[}Derived sheaves{]}. Our order
of argument starts with lifting sections, because that is the mechanism
behind the vanishing and gluing theorems. Complexes have cohomological
indexing: a differential raises degree by one. Write \(\mathcal O\) for
the sheaf of rings of a ringed space \(X\), and \(\mathcal F(U)\) for
sections on an open subset \(U\).

\subsection{1. The obstruction that cohomology
records}\label{the-obstruction-that-cohomology-records}

For an exact sequence of sheaves of modules \[
0\longrightarrow\mathcal A\longrightarrow\mathcal B
\longrightarrow\mathcal C\longrightarrow0,
\] the sequence of sections on \(U\) is exact at its first two terms.
Surjectivity at the last term need not hold. A section of
\(\mathcal C(U)\) has lifts on an open cover of \(U\), since
surjectivity of a sheaf map means surjectivity on stalks. Differences of
lifts lie in \(\mathcal A\). Changing the lifts changes those
differences by differences of sections of \(\mathcal A\). This is the
first obstruction to a global lift.

Choose an injective resolution \[
0\longrightarrow\mathcal F\longrightarrow\mathcal I^0
\xrightarrow{d^0}\mathcal I^1\xrightarrow{d^1}\cdots.
\] Define \[
H^q(U,\mathcal F)
=H^q\bigl(\Gamma(U,\mathcal I^\bullet)\bigr)
\qquad(q\geq0).
\] The derived-functor construction shows that this is independent of
the resolution, naturally in \(\mathcal F\). In degree zero it is
\(\mathcal F(U)\). Negative cohomology groups of a sheaf are zero. A
short exact sequence gives a natural long exact sequence \[
0\to\mathcal A(U)\to\mathcal B(U)\to\mathcal C(U)
\xrightarrow{\partial}H^1(U,\mathcal A)
\to H^1(U,\mathcal B)\to\cdots.
\] Thus \(\partial(c)=0\) exactly when \(c\) has a lift on \(U\). The
construction also gives restriction maps for inclusions of open sets and
their compatibility with connecting maps. These are the universal right
derived functors of \(\Gamma(U,-)\): the positive-degree functors can be
effaced by embedding a sheaf into an injective sheaf.

\textbf{Proposition 1.1 (restriction).} If \(U\subset X\) is open, then
\[
H^q(U,\mathcal F)=H^q(U,\mathcal F|_U).
\] On the left, the resolution is taken in modules on \(X\); on the
right, it is taken in modules on \(U\).

\textbf{Proof.} Restriction is exact, and the restriction of an
injective sheaf of modules is injective by
\href{https://kokunoyumeto.github.io/open-math-courses/courses/derived-categories-and-sheaf-operations/injective-modules-and-bounded-below-derived-functors.html\#1-extending-maps-and-extending-sections}{Injective
restrictions, Lemma 1.1}. Therefore \(\mathcal I^\bullet|_U\) is an
injective resolution of \(\mathcal F|_U\). The two complexes of sections
in the assertion are literally the same complex. This also proves
compatibility with further restriction. \(\square\)

The proposition is useful even when \(\mathcal F|_U\) is simpler than
\(\mathcal F\). Cohomology on \(U\) depends only on the restriction, not
on how the sheaf extends to the rest of the space. It does not say that
restriction induces an isomorphism
\(H^q(X,\mathcal F)\to H^q(U,\mathcal F)\).

\subsection{2. Extending sections removes the
obstruction}\label{extending-sections-removes-the-obstruction}

A sheaf is \textbf{flasque}, also called \textbf{flabby}, if every
restriction \[
\mathcal F(V)\longrightarrow\mathcal F(U),\qquad U\subset V,
\] is surjective. Equivalently, it suffices to test restrictions from
\(X\) to all opens: a section extended to \(X\) can then be restricted
to \(V\). A restriction of a flasque sheaf to an open subset is flasque.
The adjective concerns the underlying sheaf; it does not require any
property of the ring of coefficients.

\textbf{Lemma 2.1 (lifting with an extensible kernel).} In a short exact
sequence \(0\to\mathcal A\to\mathcal B\to\mathcal C\to0\), if
\(\mathcal A\) is flasque, then \(\mathcal B(U)\to\mathcal C(U)\) is
surjective for every open \(U\).

\textbf{Proof.} Fix \(c\in\mathcal C(U)\). Choose local lifts
\(b_\lambda\in\mathcal B(W_\lambda)\) on an open cover of \(U\), and
well order the cover. We construct a compatible lift on the union of the
sets already considered. At a successor step let \(V\) be this union
with its lift \(b\), and let \(W\) be the next open set with a local
lift \(b_W\). On \(V\cap W\), the difference \(b-b_W\) is the image of a
unique section \(a\in\mathcal A(V\cap W)\), because kernels of sheaf
maps are computed on sections. Extend \(a\) to
\(\widetilde a\in\mathcal A(W)\). Replace \(b_W\) by
\(b_W+\widetilde a\); the two lifts now agree, so they glue on
\(V\cup W\). At a limit step all previous lifts agree, and the sheaf
axiom glues them on their union. Starting with the empty open set and
proceeding through the well-ordered cover produces a lift on all of
\(U\). This uses the usual axiom of choice, with no countability
assumption on the cover. \(\square\)

\textbf{Corollary 2.2.} If both \(\mathcal A\) and \(\mathcal B\) in
Lemma 2.1 are flasque, then \(\mathcal C\) is flasque.

\textbf{Proof.} Lift a section of \(\mathcal C(U)\) to
\(\mathcal B(U)\), extend it to \(\mathcal B(V)\), and take its image in
\(\mathcal C(V)\). \(\square\)

\textbf{Theorem 2.3 (flasque acyclicity).} For a flasque sheaf of
modules \(\mathcal F\), \[
H^q(U,\mathcal F)=0\qquad(q>0)
\] on every open \(U\subset X\).

\textbf{Proof.} Injective sheaves of modules are flasque by
\href{https://kokunoyumeto.github.io/open-math-courses/courses/derived-categories-and-sheaf-operations/injective-modules-and-bounded-below-derived-functors.html\#1-extending-maps-and-extending-sections}{Injective
restrictions, Lemma 1.1}. Embed \(\mathcal F\) into an injective
\(\mathcal I^0\), and let \(\mathcal Q^0\) be the quotient. Corollary
2.2 makes \(\mathcal Q^0\) flasque. Repeat: embed \(\mathcal Q^0\) into
an injective \(\mathcal I^1\), form its flasque quotient
\(\mathcal Q^1\), and continue. This constructs an injective resolution,
with short exact sequences \[
0\to\mathcal Q^{r-1}\to\mathcal I^r\to\mathcal Q^r\to0,
\qquad\mathcal Q^{-1}=\mathcal F.
\] Lemma 2.1 makes every one of these sequences exact after taking
sections on \(U\). Consequently the complex of sections of the injective
resolution is exact in every positive degree. Its cohomology computes
the groups in question. \(\square\)

\emph{Reference:} {[}Stacks, Tag 09SY{]}. The argument applies without
change to sheaves of abelian groups, by using the constant sheaf
\(\mathbf Z\) as the ring of coefficients.

\textbf{Corollary 2.4 (flasque resolutions).} An exact resolution
\(0\to\mathcal F\to\mathcal L^0\to\mathcal L^1\to\cdots\) with all
\(\mathcal L^r\) flasque computes cohomology on every open by taking
sections.

\textbf{Proof.} Theorem 2.3 makes its terms acyclic for \(\Gamma(U,-)\).
The acyclic-resolution theorem for right derived functors,
\href{https://kokunoyumeto.github.io/open-math-courses/courses/derived-categories-and-sheaf-operations/injective-modules-and-bounded-below-derived-functors.html\#4-right-derived-functors-on-bounded-below-complexes}{Bounded-below
derived functors, Theorem 4.1}, applies. Concretely, the short exact
sequences formed with the successive kernels give dimension-shifting
isomorphisms, which identify the cohomology of this complex with that of
an injective resolution. \(\square\)

Flasqueness is sufficient for acyclicity, but acyclicity on one space is
a weaker condition. A sheaf with zero higher cohomology on \(X\) need
not allow extension between every pair of opens. This distinction
matters when constructing resolutions: flasque terms simultaneously work
on every open set.

\subsection{3. Coefficients and higher direct
images}\label{coefficients-and-higher-direct-images}

The underlying abelian sheaf remembers the addition law of a module.
Does forgetting the module structure alter cohomology? It does not,
although forgetting need not turn an injective module sheaf into an
injective abelian sheaf.

\textbf{Theorem 3.1 (independence of the coefficient category).} The
underlying groups of \(H^q(U,\mathcal F)\), computed in modules over
\(\mathcal O_X\), are naturally the cohomology groups of the underlying
abelian sheaf.

\textbf{Proof.} Take an injective resolution in modules. Forgetting the
module structure is exact: this follows on stalks from exactness of
forgetting modules to abelian groups. Every term of this resolution is
flasque, and remains flasque after forgetting. By Theorem 2.3 for
abelian sheaves, the forgotten resolution is acyclic for sections. It
therefore computes abelian-sheaf cohomology by Corollary 2.4. The
complexes of sections in the two categories have the same underlying
abelian groups and differentials. Comparison maps are natural because
both are derived from the identity on sections. \(\square\)

Now let \(f:(X,\mathcal O_X)\to(Y,\mathcal O_Y)\) be a morphism of
ringed spaces. The sheaf \(f_*\mathcal F\) is given by
\((f_*\mathcal F)(V)=\mathcal F(f^{-1}V)\); its \(\mathcal O_Y\)-action
comes from \(f^\sharp\). Define \(Rf_*\mathcal F\) using an injective
resolution and \[
R^qf_*\mathcal F=\mathcal H^q(f_*\mathcal I^\bullet).
\] Degree zero is the ordinary direct image. Higher direct images form
natural long exact sequences of sheaves.

\textbf{Theorem 3.2 (local description).} The sheaf \(R^qf_*\mathcal F\)
is the sheaf associated to the presheaf \[
V\longmapsto H^q(f^{-1}V,\mathcal F).
\] In particular, \[
(R^qf_*\mathcal F)_y
=\mathop{\mathrm{colim}}_{y\in V}H^q(f^{-1}V,\mathcal F).
\]

\textbf{Proof.} Evaluate the complex \(f_*\mathcal I^\bullet\) on \(V\).
It becomes \(\Gamma(f^{-1}V,\mathcal I^\bullet)\), whose cohomology is
the displayed group. For any complex of sheaves, its cohomology sheaf is
the sheafification of the cohomology presheaf: kernels are computed on
sections and cokernels are sheafified. Equivalently, at a stalk both
constructions give the cohomology of the stalk complex, since filtered
colimits of modules are exact. This proves the assertion, including the
formula for stalks. \(\square\)

\emph{Reference:} {[}Stacks, Tag 01E4{]}. Sheafification in this
statement is essential. There is no general assertion that sections of
\(R^qf_*\mathcal F\) on \(V\) equal \(H^q(f^{-1}V,\mathcal F)\).

\textbf{Corollary 3.3.} Higher direct images commute with restricting
the target to an open subset. They also agree, as underlying abelian
sheaves, with higher direct images computed in abelian sheaves.

\textbf{Proof.} If \(V\subset Y\) is open, the inverse images of its
open subsets are the same whether one works with \(f\) or with
\(f^{-1}V\to V\). Proposition 1.1 and Theorem 3.2 identify the
corresponding presheaves and their sheafifications. For the second
assertion apply Theorem 3.1 to each inverse image open, then sheafify.
\(\square\)

\textbf{Corollary 3.4.} A direct image of a flasque sheaf is flasque. A
flasque sheaf is \(f_*\)-acyclic: \(R^qf_*\mathcal F=0\) for \(q>0\).

\textbf{Proof.} Restrictions for \(f_*\mathcal F\) are restrictions
between inverse image opens, so are surjective. The presheaf in Theorem
3.2 is zero in positive degrees by Theorem 2.3. \(\square\)

For the map from \(X\) to a point, higher direct images are just
cohomology groups, viewed as sheaves on a point. For other maps, the
stalk formula measures cohomology over arbitrarily small neighbourhoods
of the target point. It is not in general the cohomology of the fibre;
comparison with fibre cohomology requires additional conditions, treated
later in the course.

\subsection{4. Collecting information along a
map}\label{collecting-information-along-a-map}

Taking sections on \(X\) is the composite of direct image to \(Y\) and
sections on \(Y\). The cohomology of this composite has two possible
sources: cohomology along \(f\), and cohomology of the resulting sheaves
on \(Y\).

\textbf{Theorem 4.1 (Leray spectral sequence).} For any morphism of
ringed spaces and any sheaf of modules \(\mathcal F\), there is a
natural first-quadrant spectral sequence \[
E_2^{p,q}=H^p(Y,R^qf_*\mathcal F)
\quad\Longrightarrow\quad H^{p+q}(X,\mathcal F).
\] The groups on the right are regarded as \(\mathcal O_Y(Y)\)-modules
through the map on global rings.

\textbf{Proof.} Choose an injective resolution
\(\mathcal F\to\mathcal I^\bullet\) starting in degree zero and put
\(K=f_*\mathcal I^\bullet\). Then \(K\) represents \(Rf_*\mathcal F\)
and \(\mathcal H^qK=R^qf_*\mathcal F\). Apply the full exact-couple
construction in
\href{https://kokunoyumeto.github.io/open-math-courses/courses/derived-categories-and-sheaf-operations/derived-pullback-and-pushforward.html\#6-exercises-with-checked-solutions}{Derived
pullback and pushforward, Exercise 3} to the triangles obtained from the
good upper truncations \(\tau_{\le q}K\). It constructs the pages
\(E_r\), proves \(E_{r+1}=H(E_r,d_r)\), and identifies the limiting
terms with the successive quotients of the finite filtration of
\(H^n(R\Gamma(Y,K))\). Its hypotheses are exactly a module sheaf in
degree zero and a morphism of ringed spaces; no bound on the dimensions
is needed.
\href{https://kokunoyumeto.github.io/open-math-courses/courses/derived-categories-and-sheaf-operations/derived-pullback-and-pushforward.html\#3-composition-leray-and-the-local-descriptions}{Proposition
3.1 of the same lesson} proves
\(R\Gamma(Y,Rf_*\mathcal F)\cong R\Gamma(X,\mathcal F)\), with scalars
restricted along the map of global rings. These two internal proofs
establish the terms, convergence and abutment asserted here. Truncation
triangles and their induced exact couples are natural in \(\mathcal F\),
giving naturality of the spectral sequence. \(\square\)

The differential on page \(r\) has bidegree \((r,1-r)\). For each total
degree the first-quadrant spectral sequence gives a finite filtration
whose successive quotients are the terms \(E_\infty^{p,q}\) in that
total degree. It gives a filtration, not a canonical direct-sum
decomposition.

If \(R^qf_*\mathcal F=0\) for \(q>0\), only the bottom row survives.
Hence there are natural isomorphisms \[
H^p(X,\mathcal F)\cong H^p(Y,f_*\mathcal F).
\] More generally, the first terms produce \[
0\to H^1(Y,f_*\mathcal F)\to H^1(X,\mathcal F)
\to H^0(Y,R^1f_*\mathcal F)\to H^2(Y,f_*\mathcal F)
\to H^2(X,\mathcal F).
\] The middle map records a global cohomology class locally along \(f\).
The next map is the obstruction to such local classes coming from a
global class.

\textbf{Example 4.2 (closed embeddings).} Let \(i:Z\hookrightarrow X\)
be a closed embedding of ringed spaces. For a sheaf \(\mathcal G\) on
\(Z\), the stalk of \(i_*\mathcal G\) at a point of \(Z\) is the
corresponding stalk of \(\mathcal G\), and at a point outside \(Z\) it
is zero. For the first assertion, opens of \(Z\) are intersections with
opens of \(X\); for the second, use the open complement of \(Z\). Thus
\(i_*\) is exact, by the stalk criterion for exactness. It follows that
\[
R^qi_*\mathcal G=0\ (q>0),\qquad
H^p(X,i_*\mathcal G)\cong H^p(Z,\mathcal G).
\] The last identity is the single-row case of Leray. A closed embedding
therefore transports cohomology without adding any higher direct images.

\subsection{5. Two open sets and the connecting
map}\label{two-open-sets-and-the-connecting-map}

\textbf{Theorem 5.1 (Mayer--Vietoris).} If \(X=U\cup V\) is an open
cover, there is a natural exact sequence \[
\begin{aligned}
0&\to H^0(X,\mathcal F)
\to H^0(U,\mathcal F)\oplus H^0(V,\mathcal F)
\to H^0(U\cap V,\mathcal F)\\
&\to H^1(X,\mathcal F)
\to H^1(U,\mathcal F)\oplus H^1(V,\mathcal F)
\to H^1(U\cap V,\mathcal F)\to\cdots.
\end{aligned}
\] The map from the two opens to their intersection is restriction from
\(U\) minus restriction from \(V\).

\textbf{Proof.} Use a flasque resolution \(\mathcal L^\bullet\) of
\(\mathcal F\), for example an injective resolution. In every degree the
sheaf axiom identifies the kernel of the difference map with sections on
\(X\). The difference map is onto, since any section on \(U\cap V\)
extends to \(U\) by flasqueness. We obtain a short exact sequence of
complexes \[
0\to\Gamma(X,\mathcal L^\bullet)
\to\Gamma(U,\mathcal L^\bullet)\oplus\Gamma(V,\mathcal L^\bullet)
\to\Gamma(U\cap V,\mathcal L^\bullet)\to0.
\] Its long exact cohomology sequence is the asserted sequence by
Corollary 2.4 and Proposition 1.1. Comparison maps between resolutions
make this construction natural in \(\mathcal F\); homotopic comparison
maps induce the same maps on cohomology. \(\square\)

To see the connecting map explicitly, represent a class on \(U\cap V\)
by a cocycle \(c\in\mathcal L^q(U\cap V)\). Choose
\(a\in\mathcal L^q(U)\) and \(b\in\mathcal L^q(V)\) with \(a-b=c\) on
the overlap. Since \(dc=0\), the sections \(da\) and \(db\) agree there.
They glue to a section \(h\in\mathcal L^{q+1}(X)\), and \(dh=0\). The
class of \(h\) is the connecting image. Another choice of lifts differs
by a global section in degree \(q\), so changes \(h\) by a coboundary.
Changing \(c\) by a coboundary likewise leaves the resulting class
unchanged. This explains the sign and makes the obstruction concrete.

If the positive cohomology on all three opens vanishes, the sequence
reduces to a two-term calculation. It says that \(H^1(X,\mathcal F)\) is
the cokernel of the difference map on sections and that
\(H^q(X,\mathcal F)=0\) for \(q\geq2\). The acyclicity conditions must
be checked; an arbitrary open cover does not automatically compute
cohomology from sections alone. The next lesson extends this calculation
to Čech complexes.

\subsection{6. Simple vanishing tests and a dimension
bound}\label{simple-vanishing-tests-and-a-dimension-bound}

\textbf{Example 6.1 (one-point support).} For a point \(x\in X\) and an
abelian group \(A\), the skyscraper sheaf \(i_{x*}A\) has value \(A\) on
opens containing \(x\), and zero on other opens. Every restriction is an
identity or a map onto zero. Thus the sheaf is flasque and has no
positive cohomology on any open. This does not require \(x\) to be
closed. The same example works for a module over \(\mathcal O_{X,x}\),
with the induced action on sections.

\textbf{Example 6.2 (irreducibility).} Let \(X\) be irreducible and
\(A_X\) the constant abelian sheaf. A section is a locally constant
function to \(A\) with its discrete topology. Every nonempty open in an
irreducible space is irreducible and hence connected. Such a function is
constant. Thus \(A_X(U)=A\) for nonempty \(U\), with identity
restrictions, and \(A_X(\varnothing)=0\). The sheaf is flasque.
Consequently constant coefficients have zero positive sheaf cohomology
on an irreducible space. Connectedness alone does not give this
extension property: a connected space may have disconnected opens.

The sheaves in these examples are acyclic for a direct reason,
independently of dimension. A different theorem controls all sheaves at
once. \textbf{Grothendieck's vanishing theorem}, proved in
\href{https://kokunoyumeto.github.io/open-math-courses/courses/ag-etale-cohomology/cohomological-dimension-and-the-kunneth-formula.html\#7-the-bound-for-all-finite-type-schemes}{Cohomological
dimension and the Künneth formula, Lemma 7.1}, states that if the
underlying space \(X\) is Noetherian of finite Krull dimension \(d\),
then \[
H^q(X,\mathcal F)=0\qquad(q>d)
\] for every abelian sheaf, and hence every sheaf of modules. Here
dimension is the supremum of lengths of chains of irreducible closed
subsets, not a manifold dimension. The precise reference is {[}Stacks,
Tag 02UZ{]}. It gives no vanishing assertion in degree \(d\).

For an explicit small-space example, the restriction-map description of
flabby sheaves on the Sierpiński space is worked out in {[}Derived
sheaves, Example 6 and Exercise 3{]}. It supplies a useful distinction
between a condition on a restriction map and a condition on an
individual stalk.

\subsection{7. Exercises with solutions}\label{exercises-with-solutions}

\textbf{Exercise 7.1 (easy: two point supports).} Let \(x,y\in X\),
possibly equal, and let \(A,B\) be abelian groups. Compute the positive
cohomology of \(i_{x*}A\oplus i_{y*}B\) on every open. Give its
degree-zero sections.

\textbf{Solution.} On an open \(W\) the sections are the sum of \(A\) if
\(x\in W\) and \(B\) if \(y\in W\), with each missing summand replaced
by zero. Finite sums of flasque sheaves are flasque, because each
component of a section can be extended separately. The positive groups
are therefore zero by Theorem 2.3. Even when \(x=y\), the degree-zero
value is \(A\oplus B\) on opens containing that point.

\textbf{Exercise 7.2 (easy: closed support).} For a closed embedding
\(i:Z\hookrightarrow X\) and an exact sequence of module sheaves on
\(Z\), prove exactness after applying \(i_*\). Deduce the
higher-direct-image and cohomology statements of Example 4.2.

\textbf{Solution.} At a point of \(Z\) the pushed-forward sequence is
the original stalk sequence; outside \(Z\) it is the zero sequence. Thus
it is exact at every stalk. An exact functor has zero positive right
derived functors, giving \(R^qi_*=0\). Leray has only its \(q=0\) row
and gives \(H^p(X,i_*\mathcal G)=H^p(Z,\mathcal G)\).

\textbf{Exercise 7.3 (medium: constant sheaves).} Prove the assertion
about irreducible spaces in Example 6.2, and explain precisely why the
proof does not apply to an arbitrary connected space.

\textbf{Solution.} Nonempty opens of an irreducible space intersect. Any
two nonempty opens in a nonempty open subset therefore intersect too, so
that subset is irreducible. A nonconstant locally constant function
partitions it into two disjoint nonempty opens, a contradiction. Every
section is consequently a constant element of \(A\), which extends as
the same constant to any larger open. On a connected space an open
subset can have two components; a function taking different values on
them cannot extend to a locally constant function on the whole connected
space. The proof requires irreducibility of every nonempty open.

\textbf{Exercise 7.4 (medium: a gluing obstruction).} Assume
\(X=U\cup V\) and, for all \(q>0\), \[
H^q(U,\mathcal F)=H^q(V,\mathcal F)=H^q(U\cap V,\mathcal F)=0.
\] Prove the two-term calculation after Theorem 5.1, and describe what
vanishing of the class represented by an overlap section means.

\textbf{Solution.} The long exact sequence gives \[
H^1(X,\mathcal F)=\frac{\mathcal F(U\cap V)}{\mathcal F(U)|_{U\cap V}-\mathcal F(V)|_{U\cap V}}.
\] Its later terms give zero in degrees at least two. The class of \(c\)
is zero exactly when \(c=a|_{U\cap V}-b|_{U\cap V}\). These sections
correct the mismatch of local data. The short exact sequence of
complexes in the proof establishes the calculation without assuming that
the sections functor is exact on \(\mathcal F\).

\textbf{Exercise 7.5 (medium: an acyclic direct-image test).} Suppose
every inverse image of an open subset of \(Y\) has zero positive
cohomology with coefficients in \(\mathcal F\). Show that
\(H^p(X,\mathcal F)=H^p(Y,f_*\mathcal F)\). Apply this to a closed
embedding.

\textbf{Solution.} The presheaves describing \(R^qf_*\mathcal F\) are
zero for \(q>0\), so their sheafifications are zero. Leray gives the
isomorphisms. For a closed embedding the restricted embedding into any
target open is still closed and its direct image is exact by Exercise
7.2. In this application one uses exactness to establish vanishing of
higher direct images directly; one need not assume that arbitrary opens
of \(Z\) are acyclic. This distinction is essential: the sufficient
hypothesis in the first sentence is stronger than the conclusion
\(R^qf_*\mathcal F=0\).

\textbf{Exercise 7.6 (challenging: changing rings).} Let
\(\mathcal R\to\mathcal S\) be a morphism of sheaves of rings on the
same space and \(\mathcal F\) an \(\mathcal S\)-module. Prove that its
cohomology computed as an \(\mathcal R\)-module is the restriction of
scalars of its cohomology as an \(\mathcal S\)-module.

\textbf{Solution.} Forgetting from \(\mathcal S\)-modules to
\(\mathcal R\)-modules is exact on stalks. An injective resolution over
\(\mathcal S\) becomes an exact flasque resolution over \(\mathcal R\).
Theorem 2.3 and the acyclic-resolution theorem show that this resolution
computes \(\mathcal R\)-module cohomology. Its section complex is the
restriction of scalars of the original section complex; restriction of
scalars is exact. This proves the assertion and its naturality. No
flatness of \(\mathcal S\) over \(\mathcal R\) is needed.

\subsection{Foundational cohomology
results}\label{foundational-cohomology-results}

The internal foundations are
\href{https://kokunoyumeto.github.io/open-math-courses/courses/derived-categories-and-sheaf-operations/injective-modules-and-bounded-below-derived-functors.html}{Injective
modules and bounded-below derived functors, Theorem 2.3, Lemma 1.1,
Proposition 1.3 and Theorem 4.1}: enough injectives, flasqueness and
open restriction of injectives, and the existence and acyclic-resolution
properties of right derived functors. Leray uses the full internal proof
in
\href{https://kokunoyumeto.github.io/open-math-courses/courses/derived-categories-and-sheaf-operations/derived-pullback-and-pushforward.html}{Derived
pullback and pushforward, Proposition 3.1 and Exercise 3}. The
Noetherian dimension bound uses
\href{https://kokunoyumeto.github.io/open-math-courses/courses/ag-etale-cohomology/cohomological-dimension-and-the-kunneth-formula.html\#7-the-bound-for-all-finite-type-schemes}{Cohomological
dimension and the Künneth formula, Lemma 7.1}; this particular lemma
concerns arbitrary abelian sheaves on the underlying Noetherian
topological space and has no étale torsion hypothesis. The lifting,
flasque acyclicity, coefficient comparison, local description of higher
direct images, and the applications to Leray and Mayer--Vietoris have
been proved above.

\subsection{References}\label{references}

\begin{itemize}
\tightlist
\item
  \textbf{{[}Stacks{]}} The Stacks project authors, \emph{The Stacks
  project}, Cohomology of Sheaves: Tags
  \href{https://kokunoyumeto.github.io/stacks-zh-hans-cn/en/cohomology.html\#cohomology-lemma-cohomology-of-open}{01E1},
  \href{https://kokunoyumeto.github.io/stacks-zh-hans-cn/en/cohomology.html\#cohomology-lemma-describe-higher-direct-images}{01E4},
  \href{https://kokunoyumeto.github.io/stacks-zh-hans-cn/en/cohomology.html\#cohomology-lemma-flasque-acyclic}{09SY},
  \href{https://kokunoyumeto.github.io/stacks-zh-hans-cn/en/cohomology.html\#cohomology-lemma-flasque-acyclic-pushforward}{09T0},
  \href{https://kokunoyumeto.github.io/stacks-zh-hans-cn/en/cohomology.html\#cohomology-lemma-modules-abelian}{01F1},
  \href{https://kokunoyumeto.github.io/stacks-zh-hans-cn/en/cohomology.html\#cohomology-lemma-Leray}{01F2},
  \href{https://kokunoyumeto.github.io/stacks-zh-hans-cn/en/cohomology.html\#cohomology-lemma-mayer-vietoris}{01EB},
  and
  \href{https://kokunoyumeto.github.io/stacks-zh-hans-cn/en/cohomology.html\#cohomology-proposition-vanishing-Noetherian}{02UZ};
  Derived Categories, Tag
  \href{https://kokunoyumeto.github.io/stacks-zh-hans-cn/en/derived.html\#derived-lemma-grothendieck-spectral-sequence}{015N}.
  The reference edition is the AI Integrated Stacks Project.
\item
  \textbf{{[}Derived sheaves{]}} \emph{Derived categories of sheaves},
  Open Mathematics Courses:
  \href{https://kokunoyumeto.github.io/open-math-courses/courses/derived-categories-and-sheaf-operations/injective-modules-and-bounded-below-derived-functors.html}{Injective
  modules and bounded-below derived functors} and
  \href{https://kokunoyumeto.github.io/open-math-courses/courses/derived-categories-and-sheaf-operations/derived-pullback-and-pushforward.html}{Derived
  pullback and pushforward}.
\end{itemize}
\end{document}
